\documentclass[11pt,a4paper]{article}

% --------------------------------------------------
% Packages
% --------------------------------------------------
\usepackage[utf8]{inputenc}
\usepackage[T1]{fontenc}
\usepackage[french]{babel}
\usepackage[a4paper,margin=2cm]{geometry}
\usepackage{amsmath,amssymb,amsthm,mathtools}
\usepackage{lmodern}
\usepackage{xcolor}
\usepackage{hyperref}
\usepackage{fancyhdr}
\usepackage{lastpage}
\usepackage{enumitem}
\usepackage{array}
\usepackage{tabularx}
\usepackage{stmaryrd}
\usepackage{mathrsfs}
\usepackage{tikz}
\usepackage{pgfplots}
\pgfplotsset{compat=1.18}

% --------------------------------------------------
% Hyperref
% --------------------------------------------------
\hypersetup{
    colorlinks=true,
    linkcolor=blue!50!black,
    urlcolor=blue!50!black,
    citecolor=blue!50!black
}

% --------------------------------------------------
% Header / Footer
% --------------------------------------------------
\setlength{\headheight}{14pt}
\pagestyle{fancy}
\fancyhf{}
\fancyhead[L]{Terminale générale -- Spécialité mathématiques}
\fancyhead[C]{Cours complet : Suites numériques}
\fancyhead[R]{Page \thepage/\pageref{LastPage}}
\fancyfoot[C]{Cours détaillé avec démonstrations, méthodes, exemples et ouvertures}

% --------------------------------------------------
% Theorem styles
% --------------------------------------------------
\theoremstyle{definition}
\newtheorem{definition}{Définition}[section]
\newtheorem{exemple}[definition]{Exemple}
\newtheorem{remarque}[definition]{Remarque}
\newtheorem{methode}[definition]{Méthode}
\newtheorem{vocabulaire}[definition]{Vocabulaire}
\newtheorem{interpretation}[definition]{Clé de compréhension}

\theoremstyle{plain}
\newtheorem{propriete}[definition]{Propriété}
\newtheorem{theoreme}[definition]{Théorème}
\newtheorem{proposition}[definition]{Proposition}
\newtheorem{lemme}[definition]{Lemme}
\newtheorem{corollaire}[definition]{Corollaire}

\theoremstyle{remark}
\newtheorem*{demonstrationlibre}{Démonstration}
\newtheorem*{conclusion}{Conclusion}

% --------------------------------------------------
% Commands
% --------------------------------------------------
\newcommand{\R}{\mathbb{R}}
\newcommand{\N}{\mathbb{N}}
\newcommand{\Z}{\mathbb{Z}}
\newcommand{\Q}{\mathbb{Q}}
\newcommand{\ds}{\displaystyle}
\newcommand{\e}{\mathrm{e}}
\newcommand{\eps}{\varepsilon}
\newcommand{\uu}{(u_n)_{n\in\N}}
\newcommand{\vv}{(v_n)_{n\in\N}}
\newcommand{\ww}{(w_n)_{n\in\N}}

% --------------------------------------------------
% Boxes with simple styling
% --------------------------------------------------
\newenvironment{encadre}[1]
{\begin{center}
\begin{tabular}{|p{0.93\textwidth}|}
\hline
\textbf{#1} \\
}
{\\ \hline
\end{tabular}
\end{center}}

% --------------------------------------------------
% Document
% --------------------------------------------------
\begin{document}

\begin{center}
    {\LARGE \textbf{Cours de Terminale spécialité mathématiques}}\\[0.2cm]
    {\Huge \textbf{Suites numériques}}\\[0.3cm]
    {\large Version très complète avec démonstrations, méthodes, intuition et ouvertures vers le supérieur}
\end{center}

\vspace{0.4cm}

\begin{encadre}{Objectifs du chapitre}
Dans ce chapitre, on cherche à comprendre comment décrire une évolution discrète, comment étudier son comportement à long terme, comment démontrer rigoureusement des propriétés vraies pour tout entier, et comment articuler algèbre, analyse, algorithmique et modélisation.

On veut savoir :
\begin{itemize}[leftmargin=0.8cm]
    \item définir et manipuler une suite numérique ;
    \item reconnaître les suites usuelles ;
    \item étudier le sens de variation d'une suite ;
    \item encadrer une suite et comprendre la notion de bornes ;
    \item déterminer ou conjecturer la limite d'une suite ;
    \item utiliser le raisonnement par récurrence ;
    \item étudier des suites définies explicitement ou par récurrence ;
    \item relier les suites à des phénomènes d'évolution concrets ;
    \item commencer à voir des idées qui mènent vers l'analyse du supérieur.
\end{itemize}
\end{encadre}

\tableofcontents
\newpage

% ==================================================
\section{Qu'est-ce qu'une suite ?}
% ==================================================

\subsection{Définition générale}

\begin{definition}
Une \textbf{suite numérique} est une fonction définie sur un ensemble d'entiers naturels à valeurs réelles.

Le plus souvent, on considère une fonction
\[
u:\N\to\R,\qquad n\mapsto u_n.
\]
Le réel $u_n$ est appelé le \textbf{terme de rang} $n$.
\end{definition}

\begin{remarque}
Une suite n'est rien d'autre qu'une liste ordonnée de nombres
\[
u_0,\;u_1,\;u_2,\;u_3,\dots
\]
mais cette liste est produite par une règle mathématique.
\end{remarque}

\begin{interpretation}
Une fonction ordinaire associe un nombre réel à un nombre réel.  
Une suite associe un nombre réel à un entier.  

Autrement dit, une suite décrit une évolution \emph{pas à pas}. Le paramètre n'est plus un temps continu, mais un nombre d'étapes.
\end{interpretation}

\subsection{Premiers exemples}

\begin{exemple}
La suite définie par
\[
u_n=2n+1
\]
vérifie
\[
u_0=1,\qquad u_1=3,\qquad u_2=5,\qquad u_3=7.
\]
\end{exemple}

\begin{exemple}
La suite définie par
\[
u_n=\frac1{n+1}
\]
vérifie
\[
u_0=1,\qquad u_1=\frac12,\qquad u_2=\frac13,\qquad u_3=\frac14.
\]
\end{exemple}

\begin{exemple}
La suite définie par
\[
u_n=(-1)^n
\]
alterne :
\[
u_0=1,\quad u_1=-1,\quad u_2=1,\quad u_3=-1,\dots
\]
\end{exemple}

\subsection{À partir de quel rang ?}

\begin{remarque}
Certaines suites ne sont définies qu'à partir d'un certain rang. Par exemple,
\[
u_n=\frac1n
\]
n'a de sens qu'à partir de $n=1$.

On étudie alors la suite sur $\N^*=\{1,2,3,\dots\}$.
\end{remarque}

% ==================================================
\section{Deux grands modes de définition}
% ==================================================

\subsection{Définition explicite}

\begin{definition}
Une suite est dite \textbf{définie explicitement} lorsque le terme $u_n$ est donné directement en fonction de $n$.
\end{definition}

\begin{exemple}
\[
u_n=3n^2-5n+1
\]
est une définition explicite.
\end{exemple}

\begin{remarque}
Lorsqu'une suite est donnée explicitement, on peut souvent utiliser des méthodes proches de celles de l'étude des fonctions : comparaison, signe, développement, factorisation, etc.
\end{remarque}

\subsection{Définition par récurrence}

\begin{definition}
Une suite est dite \textbf{définie par récurrence} lorsqu'on donne :
\begin{itemize}[leftmargin=0.8cm]
    \item un terme initial, par exemple $u_0$ ;
    \item une relation permettant de calculer $u_{n+1}$ à partir de $u_n$.
\end{itemize}
\end{definition}

\begin{exemple}
\[
u_0=3,\qquad u_{n+1}=u_n+2.
\]
\end{exemple}

\begin{exemple}
\[
u_0=1,\qquad u_{n+1}=\frac{u_n+4}{2}.
\]
\end{exemple}

\begin{interpretation}
Une définition explicite donne directement le terme de rang $n$.

Une définition par récurrence décrit un \emph{mécanisme d'évolution} : on sait comment passer d'un état au suivant, mais on ne connaît pas forcément immédiatement la formule du terme général.
\end{interpretation}

% ==================================================
\section{Représentation graphique d'une suite}
% ==================================================

\begin{definition}
Représenter une suite $(u_n)$ graphiquement consiste à placer dans un repère les points de coordonnées
\[
(n,u_n).
\]
\end{definition}

\begin{remarque}
Comme $n$ est entier, on n'obtient pas une courbe continue, mais un ensemble de points.
\end{remarque}

\begin{methode}
Pour tracer une suite :
\begin{enumerate}[label=\arabic*.]
    \item on calcule quelques termes ;
    \item on place les points $(n,u_n)$ ;
    \item on observe la tendance : croissance, décroissance, oscillation, stabilisation apparente.
\end{enumerate}
\end{methode}

\begin{interpretation}
Le graphique donne une intuition, jamais une preuve.  

Une suite peut sembler se stabiliser sur les premiers termes et pourtant diverger ensuite. Le dessin sert à conjecturer, pas à démontrer.
\end{interpretation}

% ==================================================
\section{Variations d'une suite}
% ==================================================

\subsection{Définitions}

\begin{definition}
Soit $\uu$.
\begin{itemize}[leftmargin=0.8cm]
    \item On dit que $(u_n)$ est \textbf{croissante} si, pour tout $n\in\N$,
    \[
    u_{n+1}\geq u_n.
    \]
    \item On dit que $(u_n)$ est \textbf{strictement croissante} si, pour tout $n\in\N$,
    \[
    u_{n+1}>u_n.
    \]
    \item On dit que $(u_n)$ est \textbf{décroissante} si, pour tout $n\in\N$,
    \[
    u_{n+1}\leq u_n.
    \]
    \item On dit que $(u_n)$ est \textbf{strictement décroissante} si, pour tout $n\in\N$,
    \[
    u_{n+1}<u_n.
    \]
\end{itemize}
\end{definition}

\begin{definition}
Une suite est dite \textbf{monotone} si elle est croissante ou décroissante.
\end{definition}

\subsection{Critère par différence}

\begin{propriete}
Pour étudier le sens de variation d'une suite, on peut étudier le signe de
\[
u_{n+1}-u_n.
\]
\begin{itemize}[leftmargin=0.8cm]
    \item Si, pour tout $n$, $u_{n+1}-u_n\geq0$, alors $(u_n)$ est croissante.
    \item Si, pour tout $n$, $u_{n+1}-u_n\leq0$, alors $(u_n)$ est décroissante.
\end{itemize}
\end{propriete}

\begin{demonstrationlibre}
Par définition,
\[
(u_n)\text{ croissante}\iff \forall n,\;u_{n+1}\geq u_n.
\]
Or
\[
u_{n+1}\geq u_n\iff u_{n+1}-u_n\geq0.
\]
De même,
\[
u_{n+1}\leq u_n\iff u_{n+1}-u_n\leq0.
\]
\end{demonstrationlibre}

\begin{exemple}
Soit
\[
u_n=n^2-4n+1.
\]
Alors
\[
u_{n+1}=(n+1)^2-4(n+1)+1=n^2-2n-2.
\]
Donc
\[
u_{n+1}-u_n=(n^2-2n-2)-(n^2-4n+1)=2n-3.
\]
Ainsi :
\begin{itemize}[leftmargin=0.8cm]
    \item pour $n=0$, on a $u_1-u_0<0$ ;
    \item pour $n=1$, on a $u_2-u_1<0$ ;
    \item pour $n\geq2$, on a $u_{n+1}-u_n>0$.
\end{itemize}
La suite est donc décroissante au début puis croissante à partir du rang $2$.
\end{exemple}

\subsection{Critère par quotient}

\begin{propriete}
Si tous les termes de la suite sont strictement positifs, on peut aussi comparer $u_{n+1}$ et $u_n$ en étudiant
\[
\frac{u_{n+1}}{u_n}.
\]
\begin{itemize}[leftmargin=0.8cm]
    \item Si $\frac{u_{n+1}}{u_n}\geq1$, alors $u_{n+1}\geq u_n$.
    \item Si $\frac{u_{n+1}}{u_n}\leq1$, alors $u_{n+1}\leq u_n$.
\end{itemize}
\end{propriete}

\begin{demonstrationlibre}
Si $u_n>0$, alors diviser par $u_n$ ne change pas le sens des inégalités :
\[
u_{n+1}\geq u_n\iff \frac{u_{n+1}}{u_n}\geq1.
\]
De même,
\[
u_{n+1}\leq u_n\iff \frac{u_{n+1}}{u_n}\leq1.
\]
\end{demonstrationlibre}

\begin{remarque}
Cette méthode est particulièrement adaptée aux suites géométriques ou à toute suite présentant naturellement un produit ou une puissance.
\end{remarque}

\subsection{Variation à partir d'un certain rang}

\begin{definition}
On dit qu'une suite est croissante \textbf{à partir d'un certain rang} s'il existe $n_0\in\N$ tel que, pour tout $n\geq n_0$,
\[
u_{n+1}\geq u_n.
\]
\end{definition}

\begin{interpretation}
En étude asymptotique, ce qui se passe au début importe peu : on s'intéresse surtout au comportement à long terme.
\end{interpretation}

% ==================================================
\section{Suites arithmétiques}
% ==================================================

\subsection{Définition}

\begin{definition}
Une suite $(u_n)$ est dite \textbf{arithmétique} s'il existe un réel $r$ tel que, pour tout $n\in\N$,
\[
u_{n+1}=u_n+r.
\]
Le réel $r$ est appelé la \textbf{raison} de la suite.
\end{definition}

\begin{exemple}
\[
u_0=5,\qquad u_{n+1}=u_n+3.
\]
C'est une suite arithmétique de raison $3$.
\end{exemple}

\subsection{Formule explicite}

\begin{theoreme}
Si $(u_n)$ est une suite arithmétique de raison $r$, alors
\[
u_n=u_0+nr.
\]
Plus généralement, pour tous entiers $p,n$,
\[
u_n=u_p+(n-p)r.
\]
\end{theoreme}

\begin{demonstrationlibre}
Montrons par récurrence que
\[
u_n=u_0+nr.
\]

\textbf{Initialisation :} pour $n=0$,
\[
u_0=u_0+0\times r.
\]
La propriété est vraie au rang $0$.

\textbf{Hérédité :} supposons qu'au rang $n$ on ait
\[
u_n=u_0+nr.
\]
Alors
\[
u_{n+1}=u_n+r=(u_0+nr)+r=u_0+(n+1)r.
\]
La propriété est donc vraie au rang $n+1$.

\textbf{Conclusion :} pour tout $n\in\N$,
\[
u_n=u_0+nr.
\]
\end{demonstrationlibre}

\subsection{Caractérisation}

\begin{propriete}
Une suite est arithmétique si et seulement si la différence
\[
u_{n+1}-u_n
\]
est constante.
\end{propriete}

\begin{demonstrationlibre}
Dire que la suite est arithmétique signifie précisément qu'il existe un réel $r$ tel que
\[
u_{n+1}=u_n+r.
\]
Cela équivaut à dire
\[
u_{n+1}-u_n=r,
\]
donc que la différence est constante.
\end{demonstrationlibre}

\subsection{Variation}

\begin{propriete}
Soit $(u_n)$ une suite arithmétique de raison $r$.
\begin{itemize}[leftmargin=0.8cm]
    \item Si $r>0$, elle est strictement croissante.
    \item Si $r<0$, elle est strictement décroissante.
    \item Si $r=0$, elle est constante.
\end{itemize}
\end{propriete}

\begin{demonstrationlibre}
On a pour tout $n$ :
\[
u_{n+1}-u_n=r.
\]
Le signe de $u_{n+1}-u_n$ est donc celui de $r$, d'où le résultat.
\end{demonstrationlibre}

\subsection{Somme des termes consécutifs}

\begin{theoreme}
Si $(u_n)$ est une suite arithmétique, alors pour tous entiers $p\leq n$,
\[
u_p+u_{p+1}+\cdots+u_n=(n-p+1)\frac{u_p+u_n}{2}.
\]
\end{theoreme}

\begin{demonstrationlibre}
Posons
\[
S=u_p+u_{p+1}+\cdots+u_n.
\]
Écrivons aussi
\[
S=u_n+u_{n-1}+\cdots+u_p.
\]
En additionnant membre à membre, on obtient :
\[
2S=(u_p+u_n)+(u_{p+1}+u_{n-1})+\cdots
\]
Or, dans une suite arithmétique, toutes les paires symétriques ont la même somme, égale à $u_p+u_n$.

Comme il y a $n-p+1$ termes,
\[
2S=(n-p+1)(u_p+u_n),
\]
donc
\[
S=(n-p+1)\frac{u_p+u_n}{2}.
\]
\end{demonstrationlibre}

\begin{exemple}
Calculer
\[
7+10+13+\cdots+49.
\]
C'est une suite arithmétique de raison $3$. Le nombre de termes vaut
\[
\frac{49-7}{3}+1=14+1=15.
\]
Ainsi
\[
S=15\times \frac{7+49}{2}=15\times 28=420.
\]
\end{exemple}

% ==================================================
\section{Suites géométriques}
% ==================================================

\subsection{Définition}

\begin{definition}
Une suite $(u_n)$ est dite \textbf{géométrique} s'il existe un réel $q$ tel que, pour tout $n\in\N$,
\[
u_{n+1}=q\,u_n.
\]
Le réel $q$ s'appelle la \textbf{raison} de la suite.
\end{definition}

\begin{exemple}
\[
u_0=2,\qquad u_{n+1}=3u_n.
\]
C'est une suite géométrique de raison $3$.
\end{exemple}

\subsection{Formule explicite}

\begin{theoreme}
Si $(u_n)$ est une suite géométrique de raison $q$, alors
\[
u_n=u_0q^n.
\]
Plus généralement, pour tous entiers $p,n$,
\[
u_n=u_pq^{n-p}.
\]
\end{theoreme}

\begin{demonstrationlibre}
Montrons par récurrence que
\[
u_n=u_0q^n.
\]

\textbf{Initialisation :} pour $n=0$,
\[
u_0=u_0q^0=u_0.
\]

\textbf{Hérédité :} supposons
\[
u_n=u_0q^n.
\]
Alors
\[
u_{n+1}=q\,u_n=q(u_0q^n)=u_0q^{n+1}.
\]

\textbf{Conclusion :} pour tout $n\in\N$,
\[
u_n=u_0q^n.
\]
\end{demonstrationlibre}

\subsection{Caractérisation}

\begin{propriete}
Une suite de termes non nuls est géométrique si et seulement si le quotient
\[
\frac{u_{n+1}}{u_n}
\]
est constant.
\end{propriete}

\begin{demonstrationlibre}
Dire que la suite est géométrique signifie qu'il existe $q$ tel que
\[
u_{n+1}=qu_n.
\]
Comme $u_n\neq0$, cela équivaut à
\[
\frac{u_{n+1}}{u_n}=q,
\]
donc à dire que ce quotient est constant.
\end{demonstrationlibre}

\subsection{Variation}

\begin{propriete}
Soit $(u_n)$ une suite géométrique de raison $q$, dont tous les termes sont strictement positifs.
\begin{itemize}[leftmargin=0.8cm]
    \item Si $q>1$, la suite est strictement croissante.
    \item Si $0<q<1$, la suite est strictement décroissante.
    \item Si $q=1$, la suite est constante.
\end{itemize}
\end{propriete}

\begin{demonstrationlibre}
Comme $u_n>0$,
\[
\frac{u_{n+1}}{u_n}=q.
\]
Donc :
\begin{itemize}[leftmargin=0.8cm]
    \item si $q>1$, alors $u_{n+1}>u_n$ ;
    \item si $0<q<1$, alors $u_{n+1}<u_n$ ;
    \item si $q=1$, alors $u_{n+1}=u_n$.
\end{itemize}
\end{demonstrationlibre}

\begin{remarque}
Si $q<0$, le signe des termes alterne en général. La suite n'est donc généralement pas monotone.
\end{remarque}

\subsection{Somme des premiers termes}

\begin{theoreme}
Si $(u_n)$ est géométrique de raison $q\neq1$, alors
\[
u_0+u_1+\cdots+u_n=u_0\frac{1-q^{n+1}}{1-q}.
\]
\end{theoreme}

\begin{demonstrationlibre}
Posons
\[
S=u_0+u_1+\cdots+u_n.
\]
Comme $u_k=u_0q^k$, on a
\[
S=u_0+u_0q+\cdots+u_0q^n.
\]
En multipliant par $q$ :
\[
qS=u_0q+u_0q^2+\cdots+u_0q^{n+1}.
\]
En soustrayant,
\[
S-qS=u_0-u_0q^{n+1}.
\]
Donc
\[
S(1-q)=u_0(1-q^{n+1}).
\]
Comme $q\neq1$,
\[
S=u_0\frac{1-q^{n+1}}{1-q}.
\]
\end{demonstrationlibre}

% ==================================================
\section{Suites bornées, majorées, minorées}
% ==================================================

\begin{definition}
Une suite $(u_n)$ est :
\begin{itemize}[leftmargin=0.8cm]
    \item \textbf{majorée} s'il existe un réel $M$ tel que, pour tout $n$,
    \[
    u_n\leq M ;
    \]
    \item \textbf{minorée} s'il existe un réel $m$ tel que, pour tout $n$,
    \[
    u_n\geq m ;
    \]
    \item \textbf{bornée} si elle est à la fois majorée et minorée.
\end{itemize}
\end{definition}

\begin{exemple}
La suite
\[
u_n=\frac{(-1)^n}{n+1}
\]
vérifie
\[
-\frac1{n+1}\leq u_n\leq \frac1{n+1}.
\]
En particulier,
\[
-1\leq u_n\leq 1.
\]
Elle est donc bornée.
\end{exemple}

\begin{interpretation}
Dire qu'une suite est bornée signifie qu'elle reste enfermée dans une bande horizontale. Elle ne peut ni monter arbitrairement haut ni descendre arbitrairement bas.
\end{interpretation}

% ==================================================
\section{Raisonnement par récurrence}
% ==================================================

\subsection{Principe}

\begin{theoreme}[Principe de récurrence]
Soit $P(n)$ une propriété dépendant d'un entier naturel $n$.

Si :
\begin{itemize}[leftmargin=0.8cm]
    \item $P(n_0)$ est vraie pour un certain entier $n_0$ ;
    \item pour tout $n\geq n_0$, l'implication
    \[
    P(n)\Longrightarrow P(n+1)
    \]
    est vraie ;
\end{itemize}
alors $P(n)$ est vraie pour tout entier $n\geq n_0$.
\end{theoreme}

\begin{interpretation}
La récurrence, c'est l'idée des dominos :
\begin{itemize}[leftmargin=0.8cm]
    \item on fait tomber le premier ;
    \item on montre que chaque domino fait tomber le suivant ;
    \item alors tous tombent.
\end{itemize}
\end{interpretation}

\subsection{Structure rigoureuse}

\begin{methode}
Toute démonstration par récurrence doit contenir :
\begin{enumerate}[label=\arabic*.]
    \item \textbf{Initialisation} : vérification au premier rang ;
    \item \textbf{Hérédité} : on suppose la propriété vraie au rang $n$ et on la montre au rang $n+1$ ;
    \item \textbf{Conclusion}.
\end{enumerate}
\end{methode}

\subsection{Exemple fondamental}

\begin{proposition}
Pour tout $n\in\N$,
\[
1+2+\cdots+n=\frac{n(n+1)}2.
\]
\end{proposition}

\begin{demonstrationlibre}
Montrons cette propriété par récurrence.

\textbf{Initialisation :} pour $n=0$,
\[
0=\frac{0\times1}{2}.
\]
La propriété est vraie.

\textbf{Hérédité :} supposons que
\[
1+2+\cdots+n=\frac{n(n+1)}2.
\]
Alors
\[
1+2+\cdots+n+(n+1)=\frac{n(n+1)}2+(n+1).
\]
On factorise :
\[
\frac{n(n+1)}2+(n+1)=(n+1)\left(\frac n2+1\right)=\frac{(n+1)(n+2)}2.
\]
La propriété est donc vraie au rang $n+1$.

\textbf{Conclusion :} pour tout $n\in\N$,
\[
1+2+\cdots+n=\frac{n(n+1)}2.
\]
\end{demonstrationlibre}

\subsection{Récurrence et suites}

\begin{remarque}
Le raisonnement par récurrence est omniprésent en théorie des suites :
\begin{itemize}[leftmargin=0.8cm]
    \item pour démontrer une formule explicite ;
    \item pour prouver un encadrement ;
    \item pour établir la positivité d'une suite ;
    \item pour montrer qu'une suite reste dans un intervalle.
\end{itemize}
\end{remarque}

% ==================================================
\section{Convergence et divergence}
% ==================================================

\subsection{Définition intuitive}

\begin{definition}
On dit que la suite $(u_n)$ \textbf{converge} vers le réel $\ell$ lorsque les termes $u_n$ se rapprochent de plus en plus de $\ell$ quand $n$ devient très grand.

On note :
\[
\lim_{n\to+\infty}u_n=\ell.
\]
\end{definition}

\begin{definition}
On dit que la suite \textbf{diverge} lorsqu'elle n'admet pas de limite réelle.
\end{definition}

\begin{exemple}
\[
u_n=\frac1{n+1}\longrightarrow 0.
\]
\end{exemple}

\begin{exemple}
\[
u_n=n\longrightarrow +\infty.
\]
\end{exemple}

\begin{exemple}
\[
u_n=(-1)^n
\]
n'a pas de limite réelle.
\end{exemple}

\subsection{Version plus rigoureuse}

\begin{definition}
On dit que $(u_n)$ converge vers $\ell\in\R$ si, pour tout $\eps>0$, il existe un entier $N$ tel que, pour tout $n\geq N$,
\[
|u_n-\ell|<\eps.
\]
\end{definition}

\begin{interpretation}
Cette définition signifie : aussi petit que soit le couloir centré sur $\ell$, à partir d'un certain rang, tous les termes de la suite entrent dans ce couloir et n'en sortent plus.
\end{interpretation}

\begin{remarque}
Au niveau Terminale, cette définition n'est pas toujours exploitée dans toute sa technicité, mais elle donne le vrai sens mathématique de la convergence.
\end{remarque}

% ==================================================
\section{Limites usuelles}
% ==================================================

\begin{propriete}
On admet les limites usuelles suivantes :
\[
\lim_{n\to+\infty}\frac1n=0,
\qquad
\lim_{n\to+\infty}\frac1{n^k}=0\quad (k\in\N^*),
\]
\[
\lim_{n\to+\infty}n=+\infty,
\qquad
\lim_{n\to+\infty}n^k=+\infty\quad (k\in\N^*).
\]
\end{propriete}

\begin{propriete}
Pour tout réel $q$ :
\begin{itemize}[leftmargin=0.8cm]
    \item si $|q|<1$, alors
    \[
    q^n\to0;
    \]
    \item si $q=1$, alors $q^n=1$ pour tout $n$ ;
    \item si $q>1$, alors
    \[
    q^n\to+\infty;
    \]
    \item si $q\leq -1$, le comportement dépend des cas et peut être oscillant ou non convergent.
\end{itemize}
\end{propriete}

% ==================================================
\section{Inégalité de Bernoulli et limite de \(q^n\)}
% ==================================================

\subsection{Inégalité de Bernoulli}

\begin{proposition}[Inégalité de Bernoulli]
Pour tout réel $x>-1$ et pour tout entier $n\in\N$,
\[
(1+x)^n\geq 1+nx.
\]
\end{proposition}

\begin{demonstrationlibre}
Montrons cette propriété par récurrence sur $n$.

\textbf{Initialisation :} pour $n=0$,
\[
(1+x)^0=1\geq1+0x=1.
\]
C'est vrai.

\textbf{Hérédité :} supposons
\[
(1+x)^n\geq1+nx.
\]
Comme $x>-1$, on a $1+x>0$, donc on peut multiplier l'inégalité par $1+x$ :
\[
(1+x)^{n+1}\geq(1+nx)(1+x).
\]
Développons :
\[
(1+nx)(1+x)=1+(n+1)x+nx^2.
\]
Or $nx^2\geq0$, donc
\[
1+(n+1)x+nx^2\geq1+(n+1)x.
\]
Ainsi
\[
(1+x)^{n+1}\geq1+(n+1)x.
\]

\textbf{Conclusion :}
\[
\forall n\in\N,\qquad (1+x)^n\geq1+nx.
\]
\end{demonstrationlibre}

\begin{interpretation}
Cette inégalité permet de comparer une puissance à une expression affine. C'est un outil très utile pour montrer qu'une suite géométrique devient très grande, ou au contraire que son inverse devient très petit.
\end{interpretation}

\subsection{Démonstration de la limite de \(q^n\) pour \(0<q<1\)}

\begin{theoreme}
Si $0<q<1$, alors
\[
\lim_{n\to+\infty}q^n=0.
\]
\end{theoreme}

\begin{demonstrationlibre}
Comme $0<q<1$, le réel
\[
\frac1q
\]
est strictement supérieur à $1$. On peut donc écrire
\[
\frac1q=1+h
\qquad \text{avec } h>0.
\]
Par l'inégalité de Bernoulli,
\[
(1+h)^n\geq1+nh.
\]
Donc
\[
\left(\frac1q\right)^n\geq1+nh.
\]
En prenant les inverses, tous les termes étant positifs,
\[
q^n\leq \frac1{1+nh}.
\]
Or
\[
\frac1{1+nh}\longrightarrow0.
\]
Par encadrement,
\[
0\leq q^n\leq \frac1{1+nh}\longrightarrow0.
\]
Donc
\[
q^n\longrightarrow0.
\]
\end{demonstrationlibre}

\subsection{Cas \(|q|<1\)}

\begin{corollaire}
Si $|q|<1$, alors
\[
q^n\longrightarrow0.
\]
\end{corollaire}

\begin{demonstrationlibre}
Comme $|q|<1$, on a
\[
|q^n|=|q|^n.
\]
D'après le théorème précédent appliqué à $|q|$,
\[
|q|^n\to0.
\]
Or
\[
-|q|^n\leq q^n\leq |q|^n.
\]
Par le théorème des gendarmes,
\[
q^n\to0.
\]
\end{demonstrationlibre}

% ==================================================
\section{Opérations sur les limites}
% ==================================================

\begin{propriete}
Si
\[
u_n\to \ell
\qquad \text{et} \qquad
v_n\to \ell',
\]
alors :
\begin{itemize}[leftmargin=0.8cm]
    \item
    \[
    u_n+v_n\to \ell+\ell' ;
    \]
    \item
    \[
    u_n-v_n\to \ell-\ell' ;
    \]
    \item
    \[
    u_nv_n\to \ell\ell' ;
    \]
    \item si $\ell'\neq0$,
    \[
    \frac{u_n}{v_n}\to \frac{\ell}{\ell'}.
    \]
\end{itemize}
\end{propriete}

\begin{remarque}
Ces règles ressemblent aux règles sur les limites de fonctions. Elles sont très puissantes, mais elles ont des conditions : on ne divise pas par une suite qui tend vers $0$ sans précaution.
\end{remarque}

\begin{exemple}
Si
\[
u_n=\frac{2n+1}{n+3},
\]
alors en divisant numérateur et dénominateur par $n$,
\[
u_n=\frac{2+\frac1n}{1+\frac3n}\longrightarrow \frac21=2.
\]
\end{exemple}

% ==================================================
\section{Comparaison de suites et théorème des gendarmes}
% ==================================================

\subsection{Comparaison terme à terme}

\begin{definition}
Dire que
\[
u_n\leq v_n \quad \text{pour tout } n
\]
signifie que chaque terme de la suite $(u_n)$ est inférieur ou égal au terme correspondant de la suite $(v_n)$.
\end{definition}

\subsection{Théorème des gendarmes}

\begin{theoreme}[Théorème des gendarmes]
Soient trois suites $(u_n)$, $(v_n)$ et $(w_n)$ telles qu'à partir d'un certain rang,
\[
u_n\leq v_n\leq w_n.
\]
Si
\[
u_n\to \ell
\qquad \text{et} \qquad
w_n\to \ell,
\]
alors
\[
v_n\to \ell.
\]
\end{theoreme}

\begin{exemple}
Étudier la limite de
\[
u_n=\frac{(-1)^n}{n}.
\]
Pour tout $n\geq1$,
\[
-\frac1n\leq \frac{(-1)^n}{n}\leq \frac1n.
\]
Or
\[
-\frac1n\to0
\qquad \text{et} \qquad
\frac1n\to0.
\]
Donc
\[
\frac{(-1)^n}{n}\to0.
\]
\end{exemple}

\begin{interpretation}
Une suite peut être compliquée à regarder directement. Si on parvient à l'enfermer entre deux suites simples ayant la même limite, alors son comportement devient clair.
\end{interpretation}

% ==================================================
\section{Suites monotones et convergence}
% ==================================================

\subsection{Résultat fondamental admis}

\begin{theoreme}[admis]
Toute suite croissante et majorée converge.

Toute suite décroissante et minorée converge.
\end{theoreme}

\begin{interpretation}
C'est l'un des théorèmes les plus importants du chapitre.  

Il dit qu'une suite qui avance toujours dans le même sens, sans pouvoir partir trop loin, finit nécessairement par se stabiliser vers une limite.
\end{interpretation}

\subsection{Cas non majoré}

\begin{proposition}
Toute suite croissante non majorée tend vers $+\infty$.
\end{proposition}

\begin{demonstrationlibre}
Dire que la suite n'est pas majorée signifie :
\[
\forall A\in\R,\;\exists n_A\in\N\text{ tel que }u_{n_A}>A.
\]
Comme la suite est croissante, pour tout $n\geq n_A$,
\[
u_n\geq u_{n_A}>A.
\]
Ainsi, à partir d'un certain rang, tous les termes dépassent $A$.

C'est exactement la définition de
\[
u_n\to+\infty.
\]
\end{demonstrationlibre}

\begin{proposition}
Toute suite décroissante non minorée tend vers $-\infty$.
\end{proposition}

\begin{demonstrationlibre}
La preuve est analogue : si la suite n'est pas minorée, alors pour tout réel $A$, on peut trouver un rang où $u_n<A$, et comme la suite est décroissante, tous les termes suivants restent $\leq A$.
\end{demonstrationlibre}

\subsection{Divergence par comparaison}

\begin{proposition}
Si, à partir d'un certain rang,
\[
u_n\geq v_n
\]
et si
\[
v_n\to+\infty,
\]
alors
\[
u_n\to+\infty.
\]
\end{proposition}

\begin{demonstrationlibre}
Soit $A\in\R$. Comme $v_n\to+\infty$, il existe $N$ tel que pour tout $n\geq N$,
\[
v_n>A.
\]
Or pour $n$ assez grand, on a aussi
\[
u_n\geq v_n.
\]
Donc, pour tout $n$ assez grand,
\[
u_n>A.
\]
C'est exactement la définition de
\[
u_n\to+\infty.
\]
\end{demonstrationlibre}

\begin{corollaire}
Si, à partir d'un certain rang,
\[
u_n\leq v_n
\]
et si
\[
u_n\to-\infty,
\]
alors
\[
v_n\to-\infty
\]
n'est pas vrai en général.

En revanche, si
\[
v_n\leq u_n
\quad \text{et}\quad
u_n\to-\infty,
\]
alors
\[
v_n\to-\infty.
\]
\end{corollaire}

\begin{remarque}
Quand on compare des limites infinies, le sens de l'inégalité compte énormément.
\end{remarque}

% ==================================================
\section{Méthodes générales pour étudier une suite}
% ==================================================

\begin{methode}[Quand la suite est explicite]
On peut :
\begin{itemize}[leftmargin=0.8cm]
    \item calculer quelques termes ;
    \item étudier le signe de $u_{n+1}-u_n$ ;
    \item comparer à une suite connue ;
    \item simplifier l'écriture pour étudier la limite.
\end{itemize}
\end{methode}

\begin{methode}[Quand la suite est récurrente]
On cherche en général à :
\begin{enumerate}[label=\arabic*.]
    \item montrer que la suite reste dans un intervalle ;
    \item montrer qu'elle est monotone ;
    \item en déduire la convergence ;
    \item calculer la limite éventuelle.
\end{enumerate}
\end{methode}

\begin{methode}[Pour déterminer une limite]
On peut utiliser :
\begin{itemize}[leftmargin=0.8cm]
    \item les limites usuelles ;
    \item les opérations sur les limites ;
    \item un encadrement ;
    \item le théorème des gendarmes ;
    \item une comparaison ;
    \item une étude de monotonie et de bornes.
\end{itemize}
\end{methode}

% ==================================================
\section{Suites définies par récurrence}
% ==================================================

\subsection{Problématique générale}

\begin{remarque}
Lorsqu'une suite est définie par récurrence, on ne possède pas forcément d'expression de $u_n$ en fonction de $n$.

Il faut alors raisonner autrement.
\end{remarque}

\subsection{Premier schéma classique}

\begin{exemple}[Étude complète]
Soit la suite définie par
\[
u_0=1,\qquad u_{n+1}=\frac{u_n+3}{2}.
\]
Étudions-la complètement.
\end{exemple}

\paragraph{1. Calcul des premiers termes}
\[
u_1=\frac{1+3}{2}=2,\qquad
u_2=\frac{2+3}{2}=\frac52,\qquad
u_3=\frac{\frac52+3}{2}=\frac{11}{4}.
\]
La suite semble croissante et semble se rapprocher de $3$.

\paragraph{2. Montrons que \(u_n\leq 3\) pour tout \(n\)}
Montrons-le par récurrence.

\textbf{Initialisation :}
\[
u_0=1\leq3.
\]

\textbf{Hérédité :} supposons $u_n\leq3$. Alors
\[
u_{n+1}=\frac{u_n+3}{2}\leq \frac{3+3}{2}=3.
\]
Donc $u_{n+1}\leq3$.

\textbf{Conclusion :}
\[
\forall n\in\N,\qquad u_n\leq3.
\]

\paragraph{3. Étudions la monotonie}
\[
u_{n+1}-u_n=\frac{u_n+3}{2}-u_n=\frac{3-u_n}{2}.
\]
Or $u_n\leq3$, donc
\[
3-u_n\geq0.
\]
Ainsi
\[
u_{n+1}-u_n\geq0.
\]
La suite est croissante.

\paragraph{4. Conclusion sur la convergence}
La suite est croissante et majorée par $3$. Elle converge donc.

\paragraph{5. Calcul de la limite}
Soit $\ell$ sa limite. Si l'on passe à la limite dans la relation
\[
u_{n+1}=\frac{u_n+3}{2},
\]
on obtient
\[
\ell=\frac{\ell+3}{2}.
\]
Donc
\[
2\ell=\ell+3,
\qquad
\ell=3.
\]

\begin{conclusion}
La suite converge vers $3$.
\[
\boxed{\lim_{n\to+\infty}u_n=3.}
\]
\end{conclusion}

\subsection{Très important : l'équation de la limite ne suffit pas}

\begin{remarque}
Quand on écrit
\[
\ell=\frac{\ell+3}{2},
\]
on détermine seulement une \textbf{valeur candidate} pour la limite.

Cela ne prouve pas encore que la suite converge.  
Pour prouver la convergence, il faut d'abord un argument indépendant : monotonie + bornes, ou autre méthode.
\end{remarque}

\begin{interpretation}
Résoudre l'équation de la limite, c'est répondre à la question :
\emph{``si la suite converge, vers quoi peut-elle converger ?''}

Ce n'est pas encore répondre à :
\emph{``est-ce qu'elle converge vraiment ?''}
\end{interpretation}

% ==================================================
\section{Suites et fonctions}
% ==================================================

\subsection{Lien fondamental}

\begin{remarque}
Une suite explicite
\[
u_n=f(n)
\]
est obtenue en évaluant une fonction $f$ aux entiers naturels.

On peut alors utiliser l'étude de $f$ pour comprendre la suite.
\end{remarque}

\begin{propriete}
Si la fonction $f$ est croissante sur un intervalle contenant tous les entiers assez grands, alors la suite définie par
\[
u_n=f(n)
\]
est croissante à partir d'un certain rang.
\end{propriete}

\begin{exemple}
La fonction
\[
f(x)=x^2-4x+1
\]
a pour dérivée
\[
f'(x)=2x-4.
\]
Elle est donc croissante sur $[2,+\infty[$. La suite
\[
u_n=n^2-4n+1
\]
est donc croissante à partir du rang $2$.
\end{exemple}

\subsection{Suites définies par \(u_{n+1}=f(u_n)\)}

\begin{remarque}
C'est une situation très importante :
\[
u_{n+1}=f(u_n).
\]
Le comportement de la suite dépend alors de la fonction $f$ :
\begin{itemize}[leftmargin=0.8cm]
    \item est-ce qu'elle envoie un intervalle dans lui-même ?
    \item est-ce qu'elle est croissante ?
    \item est-ce qu'elle rapproche les valeurs d'un point fixe ?
\end{itemize}
\end{remarque}

\begin{definition}
Un réel $\ell$ est appelé \textbf{point fixe} de $f$ s'il vérifie
\[
f(\ell)=\ell.
\]
\end{definition}

\begin{interpretation}
Quand une suite est définie par
\[
u_{n+1}=f(u_n),
\]
toute limite éventuelle doit être un point fixe de $f$.
\end{interpretation}

\subsection{Exemple de stabilité d'un intervalle}

\begin{proposition}
Soit
\[
u_{n+1}=f(u_n)
\]
avec $u_0\in I$, où $I$ est un intervalle.  
Si
\[
f(I)\subset I,
\]
alors tous les termes de la suite appartiennent à $I$.
\end{proposition}

\begin{demonstrationlibre}
Montrons par récurrence que $u_n\in I$ pour tout $n$.

\textbf{Initialisation :} $u_0\in I$ par hypothèse.

\textbf{Hérédité :} supposons $u_n\in I$. Comme $f(I)\subset I$, on a
\[
u_{n+1}=f(u_n)\in I.
\]

\textbf{Conclusion :} pour tout $n$,
\[
u_n\in I.
\]
\end{demonstrationlibre}

\begin{remarque}
Cette idée est capitale pour borner une suite récurrente.
\end{remarque}

% ==================================================
\section{Exemples complets et progressifs}
% ==================================================

\subsection{Exemple 1 : suite arithmétique}

\begin{exemple}
Soit
\[
u_0=7,\qquad u_{n+1}=u_n-4.
\]
Alors :
\begin{itemize}[leftmargin=0.8cm]
    \item c'est une suite arithmétique de raison $-4$ ;
    \item
    \[
    u_n=7-4n ;
    \]
    \item la suite est strictement décroissante ;
    \item comme $-4n\to -\infty$,
    \[
    u_n\to -\infty.
    \]
\end{itemize}
\end{exemple}

\subsection{Exemple 2 : suite géométrique}

\begin{exemple}
Soit
\[
u_0=5,\qquad u_{n+1}=0,8u_n.
\]
Alors :
\begin{itemize}[leftmargin=0.8cm]
    \item c'est une suite géométrique de raison $0,8$ ;
    \item
    \[
    u_n=5\times 0,8^n ;
    \]
    \item la suite est strictement décroissante ;
    \item comme $0<0,8<1$,
    \[
    u_n\to0.
    \]
\end{itemize}
\end{exemple}

\subsection{Exemple 3 : encadrement et gendarmes}

\begin{exemple}
Étudier la limite de
\[
u_n=\frac{\sin n}{n+1}.
\]
On sait que pour tout réel $x$,
\[
-1\leq \sin x\leq1.
\]
Donc pour tout $n$,
\[
-\frac1{n+1}\leq \frac{\sin n}{n+1}\leq \frac1{n+1}.
\]
Comme
\[
\frac1{n+1}\to0,
\]
le théorème des gendarmes donne
\[
u_n\to0.
\]
\end{exemple}

\subsection{Exemple 4 : suite récurrente non triviale}

\begin{exemple}
Soit
\[
u_0=0,\qquad u_{n+1}=\sqrt{u_n+2}.
\]
\end{exemple}

\paragraph{1. Montrons que \(0\leq u_n\leq 2\)}
Par récurrence.

\textbf{Initialisation :}
\[
u_0=0\in[0,2].
\]

\textbf{Hérédité :} supposons $u_n\in[0,2]$. Alors
\[
2\leq u_n+2\leq4.
\]
Comme la fonction racine carrée est croissante sur $[0,+\infty[$,
\[
\sqrt2\leq \sqrt{u_n+2}\leq2.
\]
Donc
\[
0\leq u_{n+1}\leq2.
\]
Ainsi
\[
u_{n+1}\in[0,2].
\]

\textbf{Conclusion :}
\[
\forall n,\qquad 0\leq u_n\leq2.
\]

\paragraph{2. Étudions la monotonie}
Comme $0\leq u_n\leq2$, on peut comparer $u_{n+1}$ et $u_n$ via l'inégalité
\[
\sqrt{u_n+2}\geq u_n.
\]
Comme les deux membres sont positifs, on peut élever au carré :
\[
u_n+2\geq u_n^2
\iff
u_n^2-u_n-2\leq0
\iff
(u_n-2)(u_n+1)\leq0.
\]
Or $u_n\in[0,2]$, donc cette inégalité est vraie. Ainsi
\[
u_{n+1}\geq u_n.
\]
La suite est croissante.

\paragraph{3. Convergence}
La suite est croissante et majorée par $2$, donc elle converge.

\paragraph{4. Limite}
Si $\ell$ est sa limite,
\[
\ell=\sqrt{\ell+2}.
\]
En élevant au carré :
\[
\ell^2=\ell+2
\iff
\ell^2-\ell-2=0
\iff
(\ell-2)(\ell+1)=0.
\]
Donc
\[
\ell=2 \quad \text{ou} \quad \ell=-1.
\]
Mais comme $u_n\in[0,2]$, la limite appartient à $[0,2]$, donc
\[
\ell=2.
\]

\begin{conclusion}
\[
\boxed{\lim_{n\to+\infty}u_n=2.}
\]
\end{conclusion}

% ==================================================
\section{Algorithmique et suites}
% ==================================================

\subsection{Calcul itératif}

\begin{methode}
Pour calculer un terme d'une suite définie par récurrence :
\begin{enumerate}[label=\arabic*.]
    \item on stocke le terme initial dans une variable ;
    \item on répète la relation de récurrence autant de fois que nécessaire.
\end{enumerate}
\end{methode}

\begin{exemple}
Pour la suite
\[
u_0=1,\qquad u_{n+1}=2u_n+3,
\]
on peut écrire en pseudo-code :
\begin{verbatim}
u <- 1
Pour i allant de 1 à n :
    u <- 2*u + 3
Fin Pour
\end{verbatim}
\end{exemple}

\subsection{Recherche d'un seuil}

\begin{methode}
On cherche parfois le plus petit entier $n$ tel que
\[
u_n\geq A
\quad \text{ou} \quad
u_n\leq A.
\]
On utilise alors une boucle \texttt{Tant que}.
\end{methode}

\begin{exemple}
Pour déterminer le plus petit $n$ tel que
\[
0,9^n<10^{-3},
\]
on peut écrire :
\begin{verbatim}
n <- 0
u <- 1
Tant que u >= 0.001 :
    u <- 0.9*u
    n <- n+1
Fin Tant que
\end{verbatim}
\end{exemple}

\subsection{Recherche d'une valeur approchée d'une constante}

\begin{remarque}
De nombreuses constantes sont approchées par des suites.  
Par exemple :
\begin{itemize}[leftmargin=0.8cm]
    \item $\sqrt2$ peut être approché par la méthode de Newton ;
    \item $e$ apparaît dans la suite $\left(1+\frac1n\right)^n$ ;
    \item certaines suites approchent $\pi$, $\ln 2$ ou le nombre d'or.
\end{itemize}
\end{remarque}

\begin{exemple}[Approximation de \(\sqrt2\)]
Considérons
\[
u_0=1,\qquad u_{n+1}=\frac12\left(u_n+\frac2{u_n}\right).
\]
Les premiers termes approchent rapidement $\sqrt2$.
\end{exemple}

\begin{remarque}
Cette suite dépasse légèrement le niveau standard de Terminale, mais elle constitue une excellente ouverture vers le supérieur.
\end{remarque}

% ==================================================
\section{Erreurs classiques à éviter}
% ==================================================

\begin{encadre}{Pièges fréquents}
\begin{itemize}[leftmargin=0.8cm]
    \item Confondre suite arithmétique et suite géométrique.
    \item Affirmer qu'une suite est croissante uniquement parce que ses premiers termes augmentent.
    \item Résoudre l'équation de la limite et croire que cela prouve la convergence.
    \item Oublier de vérifier les hypothèses avant d'utiliser un quotient.
    \item Penser qu'une suite bornée converge forcément.
    \item Penser qu'une suite monotone converge forcément.
    \item Croire qu'un graphique vaut démonstration.
\end{itemize}
\end{encadre}

\begin{remarque}
Une suite bornée peut ne pas converger : par exemple
\[
u_n=(-1)^n.
\]
Une suite monotone peut ne pas converger vers un réel : par exemple
\[
u_n=n.
\]
\end{remarque}

% ==================================================
\section{Clés de compréhension profondes}
% ==================================================

\begin{interpretation}
\textbf{1. Une suite décrit une évolution discrète.}  
On n'observe pas un mouvement continu, mais une succession d'étapes.
\end{interpretation}

\begin{interpretation}
\textbf{2. La monotonie décrit le sens du mouvement.}  
Croître, c'est aller toujours vers le haut ; décroître, c'est aller toujours vers le bas.
\end{interpretation}

\begin{interpretation}
\textbf{3. Être bornée, c'est ne pas pouvoir fuir.}  
Une suite bornée reste enfermée dans une zone finie.
\end{interpretation}

\begin{interpretation}
\textbf{4. Monotonie + bornes = stabilisation.}  
Si une suite va toujours dans le même sens sans pouvoir s'échapper, alors elle finit par tendre vers une limite.
\end{interpretation}

\begin{interpretation}
\textbf{5. Une limite décrit le comportement à long terme.}  
Deux suites peuvent être très différentes au début et pourtant avoir la même limite.
\end{interpretation}

\begin{interpretation}
\textbf{6. Une relation de récurrence est une dynamique.}  
L'expression
\[
u_{n+1}=f(u_n)
\]
signifie que l'état suivant dépend entièrement de l'état présent.
\end{interpretation}

\begin{interpretation}
\textbf{7. Une limite éventuelle d'une suite récurrente est souvent un point fixe.}
\end{interpretation}

% ==================================================
\section{Ouvertures vers le supérieur}
% ==================================================

\subsection{Premier regard sur les suites adjacentes}

\begin{definition}
Deux suites $(u_n)$ et $(v_n)$ sont dites \textbf{adjacentes} lorsque :
\begin{itemize}[leftmargin=0.8cm]
    \item $(u_n)$ est croissante ;
    \item $(v_n)$ est décroissante ;
    \item pour tout $n$, $u_n\leq v_n$ ;
    \item
    \[
    v_n-u_n\to0.
    \]
\end{itemize}
\end{definition}

\begin{propriete}
Deux suites adjacentes convergent vers la même limite.
\end{propriete}

\begin{remarque}
Cette notion est surtout étudiée dans le supérieur, mais elle prolonge naturellement les idées de bornes, d'encadrement et de convergence.
\end{remarque}

\subsection{Suites extraites}

\begin{definition}
Soit $(u_n)$ une suite. Une \textbf{suite extraite} est une suite de la forme
\[
u_{\varphi(n)},
\]
où $\varphi:\N\to\N$ est strictement croissante.
\end{definition}

\begin{exemple}
Si
\[
u_n=(-1)^n,
\]
alors la sous-suite des rangs pairs est
\[
u_{2n}=1,
\]
et la sous-suite des rangs impairs est
\[
u_{2n+1}=-1.
\]
\end{exemple}

\begin{interpretation}
Les suites extraites permettent de détecter des comportements cachés. Une suite peut ne pas converger, mais certaines sous-suites peuvent converger.
\end{interpretation}

\subsection{Notion de point fixe et itérations}

\begin{remarque}
Dans le supérieur, on étudie beaucoup les suites itératives :
\[
u_{n+1}=f(u_n).
\]
On cherche alors à savoir si les itérations convergent vers un point fixe de $f$, et à quelle vitesse.
\end{remarque}

\subsection{Vitesse de convergence}

\begin{remarque}
Au niveau Terminale, on dit surtout qu'une suite converge ou non.  
Dans le supérieur, on va plus loin :
\begin{itemize}[leftmargin=0.8cm]
    \item à quelle vitesse $u_n$ tend-il vers $\ell$ ?
    \item est-ce que l'erreur $|u_n-\ell|$ décroît comme $\frac1n$ ?
    \item comme $q^n$ ?
    \item plus vite encore ?
\end{itemize}
\end{remarque}

\begin{exemple}
Les suites géométriques avec $|q|<1$ convergent très rapidement vers $0$ :
\[
q^n\to0.
\]
À l'inverse,
\[
\frac1n\to0
\]
beaucoup plus lentement.
\end{exemple}

\subsection{Lien avec les séries}

\begin{remarque}
Dans le supérieur, après les suites, on rencontre souvent les \textbf{séries}, c'est-à-dire les sommes infinies
\[
\sum_{n=0}^{+\infty}u_n.
\]
Les suites géométriques y jouent un rôle central.
\end{remarque}

\subsection{L'idée de complétude de \(\R\)}

\begin{remarque}
Le théorème \emph{``toute suite croissante majorée converge''} repose en profondeur sur une propriété fondamentale de l'ensemble des réels : la \textbf{complétude}.

Cette idée est étudiée plus sérieusement dans le supérieur, mais on peut déjà comprendre qu'elle exprime que l'ensemble $\R$ ne possède pas de ``trous''.
\end{remarque}

\subsection{Suites de Cauchy}

\begin{definition}[ouverture]
Une suite $(u_n)$ est dite de \textbf{Cauchy} si les termes finissent par être arbitrairement proches les uns des autres :
\[
\forall \eps>0,\;\exists N,\;\forall p,q\geq N,\quad |u_p-u_q|<\eps.
\]
\end{definition}

\begin{remarque}
Dans $\R$, une suite converge si et seulement si elle est de Cauchy. C'est un résultat profond du supérieur, mais très structurant pour comprendre ce qu'est vraiment la convergence.
\end{remarque}

% ==================================================
\section{Grand bilan synthétique}
% ==================================================

\begin{center}
\renewcommand{\arraystretch}{1.6}
\begin{tabularx}{\textwidth}{|>{\raggedright\arraybackslash}p{4.4cm}|X|}
\hline
\textbf{Notion} & \textbf{À savoir absolument} \\
\hline
Suite numérique & Une suite est une fonction définie sur $\N$ ou à partir d'un certain rang. \\
\hline
Variation & Étudier $u_{n+1}-u_n$ ou parfois $\frac{u_{n+1}}{u_n}$. \\
\hline
Suite arithmétique & $u_{n+1}=u_n+r$, donc $u_n=u_0+nr$. \\
\hline
Suite géométrique & $u_{n+1}=qu_n$, donc $u_n=u_0q^n$. \\
\hline
Bornes & Majorée, minorée, bornée. \\
\hline
Convergence & $\lim u_n=\ell$ signifie que les termes se rapprochent de $\ell$. \\
\hline
Limites usuelles & $\frac1n\to0$, $n\to+\infty$, $q^n\to0$ si $|q|<1$. \\
\hline
Gendarmes & Encadrer entre deux suites de même limite. \\
\hline
Monotone + bornée & Convergence. \\
\hline
Récurrence & Outil fondamental pour démontrer une propriété vraie pour tout entier. \\
\hline
Suite récurrente & Chercher un intervalle stable, la monotonie, puis la limite éventuelle. \\
\hline
Ouverture supérieur & Suites adjacentes, suites extraites, Cauchy, points fixes, vitesse de convergence. \\
\hline
\end{tabularx}
\end{center}

% ==================================================
\section*{Conclusion générale}
% ==================================================

Le chapitre des suites est fondamental, car il introduit simultanément :
\begin{itemize}[leftmargin=0.8cm]
    \item une nouvelle manière de modéliser des phénomènes ;
    \item un premier accès rigoureux à l'idée de limite ;
    \item le raisonnement par récurrence ;
    \item les premières vraies techniques d'analyse ;
    \item des idées profondes qui seront prolongées dans le supérieur.
\end{itemize}

On y apprend à distinguer :
\begin{itemize}[leftmargin=0.8cm]
    \item le comportement local et le comportement à long terme ;
    \item l'intuition graphique et la démonstration ;
    \item la recherche d'une limite candidate et la preuve de convergence ;
    \item les suites explicites et les suites dynamiques définies par récurrence.
\end{itemize}

C'est donc un chapitre central, à la fois formateur, technique et très riche conceptuellement.

\end{document}